\section{Post-training I：SFT 与数据工程}

\subsection{SFT 的角色：把潜在能力变成可调用行为}
SFT 不是 ``补课''，而是建立交互协议。它让底座模型知道任务边界、输出风格和回答结构，直接影响后续偏好优化与 RL 的稳定性。

\begin{importantbox}{讲者给出的实务判断}
若 SFT 数据结构混乱、指令覆盖窄、质量参差不齐，后续 DPO/PPO 与 RLVR 会更容易出现 reward hacking、模式坍缩或能力偏科。
\end{importantbox}

\begin{figure}[H]
\centering
\includegraphics[width=0.9\textwidth]{frame-04-sft-data.jpg}
\caption{视频约 00:29:30：SFT 阶段的数据来源与清洗策略，强调人类数据与合成数据的协同}
\end{figure}

\subsection{人类数据与合成数据的混合策略}
\begin{knowledgebox}{混合的目标不是 ``谁替代谁''，而是 ``谁补谁''}
\begin{itemize}
    \item 人类标注数据：高质量、风格真实，但规模小且成本高；
    \item 合成数据：规模大、覆盖广，但质量波动与偏差风险更高；
    \item 最优策略通常是高质量人类样本定锚，合成样本扩展覆盖。
\end{itemize}
\end{knowledgebox}

讲者多次提到数据去污染（decontamination）与许可证治理，这对开放社区尤为关键。因为在 benchmark 污染存在时，训练改进会被误读为算法进步。

\begin{warningbox}{SFT 最常见的误区}
只看总量不看分布。堆更多 instruction 数据若破坏任务平衡，可能导致推理能力上升但事实性、格式遵循或多语言能力下降。
\end{warningbox}

\subsection{本章小结}
SFT 阶段的核心是数据工程而不是单一算法。可复现的数据治理流程，是后续偏好优化与强化学习稳定工作的前置条件。


